\documentclass[11pt]{article}
\usepackage[margin=1in]{geometry}
\usepackage{amsmath,amssymb,amsthm,mathtools}
\usepackage[hidelinks]{hyperref}
\newtheorem{theorem}{Theorem}
\newtheorem{lemma}[theorem]{Lemma}
\newtheorem{corollary}[theorem]{Corollary}
\newcommand{\Q}{\mathbb Q}
\newcommand{\C}{\mathbb C}
\newcommand{\F}{\mathbb F}
\newcommand{\Z}{\mathbb Z}
\newcommand{\vp}{v_p}
\title{A prime-power two-term Strong Factorial theorem via Newton residuals}
\author{Internal written proof and exact certificates; no priority claim}
\date{16 November 2026}
\begin{document}
\maketitle
\noindent\textbf{Result.}
For every prime $p$, $a\ge1$, $d=p^a$ and complex $A,B$, the nonzero
polynomial $A+BY^d$ has a nonzero factorial moment in every pair of
consecutive positive exponents. More generally this holds for the
$\operatorname{Gamma}(\alpha,1)$ functional with integer $1\le\alpha\le d$.
Thus the two-term SFC case is proved for every prime-power exponent,
uniformly in the starting index. The proof is not an irreducibility-density
argument and does not infer an all-index result from a finite test.

\textbf{Credit and scope.}
The factorial-unit congruence and Newton-polygon vertices are in
\href{sources/rocks-2015-author-preprint.pdf}{Rocks's 2015 author preprint,
Lemma 3.8 and Proposition 3.22}. That source explicitly states the
constant-plus-square case in Theorem 3.24. Here the residual polynomial
on \emph{each} edge is computed, and its squarefreeness is used with the
adjacent-moment derivative identity proved in Section 4 below.
The source's full final 2017 journal text remains unavailable; no novelty
or priority claim is made. This is an internal proof, not outside review
or formal verification. Other exponents in the ordinary constant-plus-power family are covered by the
\href{ALL_UNIVARIATE_TWO_TERM_SFC.tex}{all-index theorem}; unrestricted SFC remains open.

\section{Why residual squarefreeness excludes a square factor}
Normalize the valuation on $\Q_p$ by $\vp(p)=1$.
For $P(X)=\sum c_kX^k$ with $c_0c_n\ne0$, use the lower Newton polygon
with coordinates $(k,\vp(c_k))$. Let an edge run from $(L,V)$ with slope
$h/e$ in lowest terms, $e>0$, and horizontal length $me$.
Its normalized residual polynomial is
\[
 R(T)=\sum_{j=0}^{m}
 \overline{\frac{c_{L+je}}{p^{V+jh}}}\,
 \overline{\frac{c_L}{p^V}}^{-1}T^j\in\F_p[T].       \tag{1}
\]
Terms strictly above the edge contribute zero; the two end coefficients
are nonzero. All coefficients of factors in $\Q_p[X]$ have integer
valuations, which is important in the following argument.

\begin{lemma}\label{lem:residual}
If every edge residual (1) is squarefree, then $P$ is squarefree in
$\Q_p[X]$.
\end{lemma}
\begin{proof}
For a rational slope $\lambda=h/e$ define the weighted minimum
$w_\lambda(H)=\min_k(\vp(h_k)-\lambda k)$.
The weighted initial form of a product is the product of the weighted
initial forms. One can verify this in a valued extension with an element
$\pi$ satisfying $\pi^e=p$: substitute $X=\pi^{-h}Z$, normalize by the
smallest coefficient valuation, and reduce. The nonzero initial
polynomials multiply over the residue field, an integral domain, so their
product does not vanish. This proves both multiplicativity and the usual
addition of Newton edges, without a factor-lifting assumption.

For a polynomial over $\Q_p$, two indices attaining this same weighted
minimum differ by a multiple of $e$, because $\gcd(h,e)=1$ and their
coefficient valuations are integers. After removing its lowest monomial,
its initial form therefore has the form $R_H(Z^e)$, up to a nonzero scalar,
and (1) is exactly the polynomial $R_H(T)$, normalized at its constant.
Residual polynomials consequently multiply, up to nonzero constants, at
any fixed slope. A nontrivial edge gives a positive-degree residual with
nonzero constant coefficient.

If $P=H^2J$ for a nonconstant $H\in\Q_p[X]$, then $H(0)\ne0$.
Choose any edge of $H$. Multiplicativity makes its slope an edge slope of
$P$, and its positive-degree residual occurs squared in the residual of
$P$ at that slope. This contradicts the hypothesis. Hence $P$ has no
nonconstant square divisor.
\end{proof}

This lemma uses residual polynomials, not just distinct edge slopes.
Repeated factors on one edge are precisely why a list of slopes alone is
insufficient. No separability of $Z^e$ in characteristic $p$ is assumed;
the polynomial being tested is in the variable $T$, before replacing it
by $Z^e$. In our application $e$ is often divisible by $p$.

\section{A falling-factorial unit deformation theorem}
For an integer $n\ge1$ put $n^{\underline k}=n!/(n-k)!$.
Suppose $u_0,\ldots,u_n\in\Z_p^\times$ and
\[
 u_k/u_0\equiv\gamma^k\pmod p
 \qquad(\gamma\in\F_p^\times).
\]
Define $P_n(X)=\sum_{k=0}^n n^{\underline k}u_kX^k$.
Write $n=\sum_i a_i p^i$, with $0\le a_i<p$; set
$s_i=\sum_{j\le i}a_jp^j$ and $s_{-1}=0$.
Only indices $i$ with $a_i>0$ give edges. Put
\[
 e_i=p^i,\qquad h_i=(p^i-1)/(p-1).
\]

\begin{lemma}\label{lem:edges}
The edge for digit $a_i$ has horizontal length $a_ie_i$, vertical length
$a_ih_i$ and slope $h_i/e_i$. Its residual polynomial is
\[
 R_i(T)=\sum_{j=0}^{a_i}
       \frac{a_i!}{(a_i-j)!}(\epsilon_iT)^j,
 \qquad \epsilon_i=(-1)^{h_i}\gamma^{e_i}\in\F_p^\times. \tag{2}
\]
\end{lemma}
\begin{proof}
The units do not affect valuations, so
$\vp(c_k)=\vp(n!)-\vp((n-k)!)$.
Legendre's formula at the vertices $s_i$ gives successive slopes
$h_i/e_i$, which strictly increase with $i$.
For $s_{i-1}<k<s_i$, put $\delta=s_i-k$.
The base-$p$ digits of $n-s_i$ and $\delta$ occupy disjoint places, whence
\[
 \vp(c_k)-\left\{\vp(c_{s_i})-\delta h_i/e_i\right\}
 =\frac{\operatorname{digitsum}_p(\delta)-\delta/p^i}{p-1}\ge0.
\]
Equality holds exactly when $\delta$ is a multiple of $p^i$, since every
lower digit has weight strictly less than one after division by $p^i$.
Thus the stated segments really are the lower hull, with exactly the
edge lattice indices $k=s_{i-1}+jp^i$, $0\le j\le a_i$.

It remains to compute their units. Let
$U_p(m!)=p^{-\vp(m!)}m!$. Splitting off multiples of $p$ and applying
Wilson at each step gives, for base-$p$ digits $m_\ell$,
\[
 \overline{U_p(m!)}=(-1)^{\vp(m!)}\prod_\ell m_\ell!\in\F_p^\times.
\]
Set $N=n-s_{i-1}$. For the ratio of the falling-factorial units at the
edge indices, only digit $a_i$ changes, so
\[
 \frac{U_p(N!)}{U_p((N-jp^i)!)}
   \equiv(-1)^{jh_i}\frac{a_i!}{(a_i-j)!}\pmod p.
\]
The ratio $u_{s_{i-1}+jp^i}/u_{s_{i-1}}$ is $\gamma^{jp^i}$ modulo $p$.
Combining these two ratios proves (2), with the normalization in (1).
\end{proof}

\begin{theorem}\label{thm:deformation}
Every $P_n$ in this unit deformation class is squarefree over $\Q_p$.
\end{theorem}
\begin{proof}
For $0\le m<p$, the truncated exponential
$E_m(Z)=\sum_{j=0}^m Z^j/j!$ is squarefree over $\F_p$.
For $m\ge1$, $E_m'=E_{m-1}$ and
$E_m-E_{m-1}=Z^m/m!$; a common root would be zero, but $E_m(0)=1$.
The case $m=0$ is constant. The residual (2) is the reversal
\[
 R_i(T)=a_i!(\epsilon_iT)^{a_i}
                 E_{a_i}((\epsilon_iT)^{-1}).
\]
It has nonzero constant, and inversion and scaling preserve its simple
nonzero roots. All residuals are therefore squarefree, and
Lemma~\ref{lem:residual} applies.
\end{proof}

This is a theorem over $\Q_p$, whose coefficient valuation group is
$\Z$. The statement is not silently extended to arbitrary ramified
coefficient fields. Our factorial, Gamma and marker applications below
have rational parameter-polynomial coefficients, so this field boundary
causes no restriction on the eventual complex phase parameters.

\section{Prime-power factorial and Gamma polynomials}
For $d=p^a$, set $b=(d-1)/(p-1)$. The factorial splitting used in the
preceding audit proves
\[
 U_k=\frac{(dk)!}{p^{bk}k!}\in\Z_p^\times,
 \qquad U_k\equiv(-1)^{bk}\pmod p.                 \tag{3}
\]
Consequently
\[
 f_{n,d}(x)=\sum_{k=0}^n\binom nk(dk)!x^k,
 \qquad f_{n,d}(X/p^b)=
       \sum_{k=0}^n n^{\underline k}U_kX^k
\]
is in Theorem~\ref{thm:deformation}, with $\gamma=(-1)^b$.
It is squarefree for \emph{every} $n$, not just a set of indices of
positive density.

For the Gamma shape $\alpha$, let
\[
 f_{n,d,\alpha}(x)=\sum_{k=0}^n\binom nk(\alpha)_{dk}x^k.
\]
If $1\le\alpha\le d$ is an integer, then
\[
 \frac{(\alpha)_{dk}}{(dk)!}
 =\prod_{j=1}^{\alpha-1}(1+dk/j)\equiv1\pmod p.      \tag{4}
\]
Each factor is a $p$-adic unit: $j<d=p^a$ implies $\vp(j)<a$.
The product is interpreted in $\Z_p$, with empty product one.
Multiplying $U_k$ by (4) preserves the same geometric unit residues.
Thus $f_{n,d,\alpha}$ is also squarefree over $\Q_p$, hence over $\Q$
and $\C$. Here ``over $\C$'' refers to the rational polynomial's root
multiplicities after extending constants.

\section{The exact two-consecutive-moment consequence}
For any sequence $m_k$, the binomial polynomials
$F_n(x)=\sum_{k=0}^n\binom nk m_kx^k$ satisfy
\[
 nF_{n-1}=nF_n-xF_n'.                              \tag{5}
\]
For $m_k=(\alpha)_{dk}$, these are precisely the moments
$f_{n,d,\alpha}(x)=\mathbb E[(1+xY^d)^n]$ with
$Y\sim\operatorname{Gamma}(\alpha,1)$.
Since $f_{n,d,\alpha}(0)=1$, any common root of adjacent polynomials
would by (5) be a multiple root of the higher one. Conversely every
multiple root is common with its predecessor. The established
squarefreeness therefore proves:

\begin{theorem}\label{thm:sfc}
For every prime power $d=p^a\ge2$, integer $1\le\alpha\le d$, every
complex $A,B$ with $A+BY^d\ne0$, and every $n\ge1$,
\[
 \bigl(\mathbb E[(A+BY^d)^n],
       \mathbb E[(A+BY^d)^{n+1}]\bigr)\ne(0,0).
\]
For $\alpha=1$ this is Strong Factorial for the whole constant-plus-
prime-power-monomial family. If one coefficient is zero its moments are
directly all nonzero; otherwise factor out $A^n$ and use $x=B/A$.
\end{theorem}

By the checked independent-linear adjunction theorem, adding $t$ new
nonzero linear variable terms to a two-term member gives a window of
length $t+2$, equal to its factorial monomial count. The coefficients may
be arbitrary complex numbers. This produces all-dimensional SFC classes
with one prime-power term and otherwise constant and independent linear
terms; no powers of this entire inhomogeneous phase are inferred to
have their own monomial-count windows.

\section{Low-degree marked moments}
Write a rational marker in divided powers as
\[
 M(Y)=\sum_{e=0}^{d-1}a_eY^e/e!,\qquad C=\sum_e a_e=\mathbb E[M(Y)].
\]
Assume $C\ne0$ and each $a_e/C$ is $p$-integral. For $Y$ exponential,
\[
 \frac{\mathbb E[M(Y)Y^{dk}]}{C(dk)!}
 =\sum_e\frac{a_e}{C}\prod_{j=1}^{e}(1+dk/j)
 \equiv1\pmod p.                                  \tag{6}
\]
The same unit deformation theorem applies to
$C^{-1}\mathbb E[M(Y)(1+xY^d)^n]$. Its constant is one and it is
squarefree for every $n$. Identity (5), valid for any fixed marker,
then proves a two-consecutive-moment bound for these marked moments
of every nonzero $A+BY^d$. In particular every monomial marker $Y^e$
with $0\le e<d$ is allowed after harmless scaling by $e!$; this also
explains the Gamma shape range above.
If $A=0$, (6) directly makes each positive marked moment nonzero; if
$B=0$, the nonzero marker mass $C$ does the same. Thus neither degenerate
coefficient case requires division by a zero phase parameter.

The condition on the rational divided-power coefficients is sufficient,
not a decision outside the class. A zero-mass marker may have an
identically zero sequence at constant phase, so it cannot be included in
this theorem without a separate branch. Markers of degree at least $d$
and arbitrary combinations with $p$-adic cancellation are not covered.
The Gaussian bridge gives the corresponding even-block radial phases
and these specified averaged markers. No all-marker or universal
Gaussian-purity assertion is added.

\section{The two moment equations retain ordinary nilpotent structure}
Let $d,\alpha$ be in Theorem~\ref{thm:sfc}, and define the universal
homogeneous moment polynomial
\[
 \mu_k(A,B)=\sum_{j=0}^k\binom kj(\alpha)_{dj}A^{k-j}B^j.
\]
For a fixed $n\ge1$, put
$I_n=(\mu_n,\mu_{n+1})\subset\Q[A,B]$ and $J=(A,B)$.
The theorem shows that the two generators have no common zero in
$\mathbb P^1(\C)$, hence no common polynomial factor. In the UFD
$\Q[A,B]$ they form a regular sequence of degrees $n,n+1$.
The exact graded multiplication sequences give the Hilbert series
\[
 \operatorname{Hilb}_{\Q[A,B]/I_n}(t)
  =\frac{(1-t^n)(1-t^{n+1})}{(1-t)^2}
  =(1+t+\cdots+t^{n-1})(1+t+\cdots+t^n).          \tag{7}
\]
It follows that this is a local Artinian complete intersection of length
$n(n+1)$, with highest nonzero degree $2n-1$ of dimension one. In particular
\[
 J^{2n}\subset I_n,\qquad J^{2n-1}\not\subset I_n. \tag{8}
\]

\begin{corollary}\label{cor:ordinary}
Over every commutative $\Q$-algebra $R$, two vanishing consecutive moments
$\mu_n(A,B)=\mu_{n+1}(A,B)=0$ force $(A,B)^{2n}=0$.
Thus $(A+BY^d)^{2n}=0$ as a polynomial, and every original marked moment
vanishes at every exponent at least $2n$. This universal exponent is sharp.
\end{corollary}
\begin{proof}
Specialize the polynomial containments (8). For sharpness use
$R=\Q[A,B]/I_n$. In that ring $J^{2n-1}\ne0$. The coefficients of
$(A+BY^d)^m$ are the distinct monomials
$\binom mj A^{m-j}B^j$; the binomial coefficients are units in every
$\Q$-algebra. They generate exactly $J^m$. Hence the phase's
$(2n-1)$st power is nonzero in the universal quotient.
\end{proof}

This is ordinary, potentially nonreduced scheme data, not merely the
radical assertion $\sqrt{I_n}=J$. Over rings the two equations need not
force $A=B=0$. The even-block Gaussian substitution retains the actual
coefficient ideal $J$: the constant coefficient is $A$, and a pure
coordinate term of the radial $d$th power has coefficient a nonzero
rational multiple of $B$. Thus the same pair ideal and sharp nilpotence
bound apply to the corresponding Gaussian coefficient families.

\section{Exact certificates and open cases}
\href{prime_power_windows.py}{The exporter} recognizes a prime-power
exponent and the stated shape/marker range. For any supplied $n$ it also
constructs the actual moment polynomial, scales its parameter, builds
its Newton hull from the coefficient valuations, and compares every
literal residual against (2). It exports an all-index class certificate
only under the written theorem's hypotheses. Unsupported exponents and
shapes are rejected, and failed marker criteria make no zero decision.
The same certificate records the consecutive-pair complete-intersection
degrees, length and sharp coefficient nilpotence bound from (7)--(8).

Controls include independent rational gcd/B\'ezout identities, direct
Gamma moments, generic geometric-unit deformations, explicit repeated
factor countercontrols, digit vertices, edge lattice points and residual
coefficients, truncated exponential separability, marker conditions and
Gaussian/disjoint-variable translations. A finite collection checks the
implementation; all $n,p,a$ quantifiers use the written proof above.
Finite Groebner controls also check ordinary quotient lengths, all degree-
$2n$ monomial containments and a surviving degree-$2n-1$ coefficient.

For a non-prime-power exponent such as $d=6$, our coefficient-unit
argument has not supplied an analogous theorem. General two-term SFC,
unrestricted SFC, higher-dimensional FC outside the closed classes,
remaining EIT V/VI audits and the complete primary bibliography remain
open. The full original goal stays active.
\end{document}
